\documentclass{article}

\usepackage{preamble}
\usepackage{mathtools}
\usepackage{kpfonts}

\title{Interval Matching on the Circle}

\begin{document}
\maketitle{}

We assuming $(L_i, R_i) \in [0, 2\pi)$ and always consider the counterclockwise arc.

Let $\{L_i\}$ be the set of all left endpoints. Note that $L_i \in [0, 2\pi)$. For each $L_i$, we unenroll the circle onto the line:
\begin{align*}
  R_j \xmapsto{\alpha} \begin{cases}
    R_i &\text{if } R_i < L_j \text{ and } L_i < R_j,\\
    R_j &\text{otherwise.}
  \end{cases}
\end{align*}

Run the regular algorithm on each of these instances; if any of these instances finds a matching then output the matching, otherwise indicate none exist.

We prove correctness now. It suffices to show that whenever such a matching exists, our algorithm also finds a matching. Let $\{P_i\}$ denote this matching with $P_i \in [0, 2\pi]^2 \times [0, 2\pi]$. Fix a cyclic ordering $P_1 \prec \hdots \prec P_n$.  

Perform our transformation using $P_1 = ((L_1, R_1), x_1)$. It suffices to argue that $\{P_i\}$ corresponds to a matching $\{\alpha(P_i)\}$. By correctness of our regular algorithm, we know that we must identify a matching too then. 

\end{document}
